\section{The delay test, and what the archimedean expansion selects}
\label{sec:delay}

Condition~\ref{cond:bilinear} is the one test in this paper a theory can fail.
This section runs it. The answer is that the test as it was first posed is failed
by every correlator and passed by a free Hamiltonian, so it discriminates nothing;
but the \emph{expansion} of the delay does discriminate, one coefficient at a
time, and it recovers the arithmetic invariants of an $L$-function in order.

\subsection{The object the test can be run on}

The natural reading of Condition~\ref{cond:bilinear} is to take a protected
two-point function on the open Wilson line of \cite{OWL} and extract the phase of
its spectral kernel. There is no such correlator: \cite{OWL} establishes the
existence of supersymmetric open Wilson lines --- only in the superconformal
versions of the $\mathcal N=2$ theories, and in $\mathcal N=4$ with a defect
hypermultiplet, where the configuration is a semicircle conformally related to the
ray --- and computes \emph{expectation values}. It contains no two-point function
of open Wilson lines and no correlator of insertions along the line.\footnote{Its
abstract was read directly; the body through a secondary reading, so this is
"not found" rather than "certainly absent".}

The test is therefore run on the object that does exist and is known exactly, the
defect CFT on the $\frac12$-BPS Wilson line, whose protected two-point functions
are $C_\Delta(\lambda)/|t_{12}|^{2\Delta}$ with $\Delta=1$ for the five scalars
\cite{GRT}, $\frac32$ for the fermions, $2$ for the displacement
\cite{CHMS}, and $L$ for the $(Y\!\cdot\!\Phi)^L$ tower \cite{GiombiKomatsu} ---
every dimension protected, the coupling dependence entirely in the normalization.

\subsection{The delay of a defect two-point function}

In the ray coordinate of Section~\ref{subsec:conformal}, $x=\log r$ and
$u=x_1-x_2$, a defect two-point function is
$G_\Delta(u)=(2\sinh\frac u2)^{-2\Delta}$, the conformally covariant form of
$C/|r_1-r_2|^{2\Delta}$; the normalization is a positive constant and drops out of
every phase. Its causal symbol is
\begin{equation}\label{eq:defectsymbol}
 \widehat k_+(\tau)=\int_0^\infty e^{-i\tau u}G_\Delta(u)\dd u
 =\Gamma(1-2\Delta)\,\frac{\Gamma(\Delta+i\tau)}{\Gamma(1-\Delta+i\tau)} ,
\end{equation}
checked against quadrature. Using
$\frac{d}{d\tau}\arg\Gamma(a+i\tau)=\re\psi(a+i\tau)$ and
$\psi(1-z)-\psi(z)=\pi\cot\pi z$:

\begin{proposition}[The delay of a defect primary]\label{prop:defectdelay}
For every $\Delta$,
\[
 \mathcal T_\Delta(\tau)=-\frac{d}{d\tau}\arg\widehat k_+(\tau)
 =\pi\re\cot\big(\pi(\Delta+i\tau)\big)
 =\frac{\pi\sin(2\pi\Delta)}{\cosh(2\pi\tau)-\cos(2\pi\Delta)} .
\]
It decays like $2\pi\sin(2\pi\Delta)e^{-2\pi\tau}$; for $0<\Delta<\frac12$ its
total is $\int_0^\infty\mathcal T_\Delta=\pi(\frac12-\Delta)$; and it is periodic
in $\Delta$ with period one.
\end{proposition}

\begin{corollary}[The protected sector gives exactly nothing]\label{cor:protected}
$\mathcal T_\Delta\equiv0$ whenever $2\Delta\in\Z$, hence for every protected
operator of the $\frac12$-BPS Wilson-line defect CFT. Not small: zero, at every
frequency and every coupling.
\end{corollary}

The target grows: $2\theta'(\tau)=\log\frac\tau{2\pi}+O(\tau^{-2})$. At $\tau=3$
and $\Delta=\frac14$ the delay is $4.09\times10^{-8}$; by $\tau=10$ it is
$3\times10^{-27}$ and the target is $+0.464$.

\subsection{No correlator can wind}

The zero of Corollary~\ref{cor:protected} is an accident of half-integer
dimensions. The following is not.

Radial quantization of the ray about the pinned endpoint gives the
boundary-channel decomposition
$G(u)=\sum_n|\langle\Omega|O|n\rangle|^2e^{-\Delta_nu}=\int_{[0,\infty)}e^{-\Delta
u}\dd\mu(\Delta)$ with $\mu\geq0$, using only that the defect Hilbert space has a
positive inner product and that the dilatation generator is bounded below. For
$G_\Delta$ this is the primary with its descendants,
$\mu=\sum_k\frac{\Gamma(2\Delta+k)}{k!\Gamma(2\Delta)}\delta_{\Delta+k}$, with
manifestly positive coefficients. By Bernstein's theorem $G$ is then completely
monotone and $\widehat k_+$ is a Stieltjes transform.

\begin{proposition}[No winding]\label{prop:nowinding}
Let $G$ be the ray kernel of a two-point function in a unitary defect theory whose
dilatation generator is bounded below. Then $\re\widehat k_+>0$ on the whole real
axis --- so $\widehat k_+$ has no zeros there --- and for $\tau>0$,
$\arg\widehat k_+(\tau)\in(-\frac\pi2,0)$, whence
\[
 \int_0^T\mathcal T(\tau)\dd\tau<\frac\pi2\qquad\text{for every }T .
\]
If a real subtraction is needed for convergence the bound is $\pi$. Either way it
is a constant, while $\int_0^T2\theta'=2\theta(T)$ is $175.9$ at $T=100$ and
$4069.1$ at $T=10^3$.
\end{proposition}

\begin{proof}
$\widehat k_+(\tau)=\int(\Delta+i\tau)^{-1}\dd\mu$ has
$\re\widehat k_+=\int\Delta(\Delta^2+\tau^2)^{-1}\dd\mu>0$ and
$\im\widehat k_+=-\tau\int(\Delta^2+\tau^2)^{-1}\dd\mu<0$. Then
$\int_0^T\mathcal T=\arg\widehat k_+(0^+)-\arg\widehat k_+(T)$.
\end{proof}

Two readings. \textbf{Phase winding is zeros crossing the contour}, and a positive
spectral measure forbids them: the scattering dictionary of
Section~\ref{sec:allpass} lists the zeros of $\xi$ as the resonances of
$K_\omega$, and a correlator of a unitary theory has none to list. And
\textbf{the obstruction is that the generator is bounded below}, not that its
spectrum is discrete --- a continuum on $[0,\infty)$ gives the same Stieltjes
transform.

\begin{remark}[The three escapes, one of which is already a trap]\label{rem:escapes}
Given those hypotheses, winding requires exactly one of: a signed measure, that
is a non-unitary theory; a spectrum \emph{not} bounded below, so that the
imaginary axis lies outside the region of absolute convergence and the object on
it is a continuation past poles --- which is precisely
Remark~\ref{rem:trap}, seen from the candidate's side rather than the target's;
or a \emph{ratio}. The Weil family takes the third route, and
Proposition~\ref{prop:groupdelay} says so: the same real function
$b(\tau^2)+w_0$ is the archimedean multiplier of the form, where it is a positive
weight with no phase, and the phase derivative of the transfer, where it winds. A
correlator offers only the first reading.
\end{remark}

\subsection{What does wind, and why that is a problem}\label{ssec:winds}

The phase shift of the inverted harmonic oscillator on the half-line is
$\eta(\tau)=\arg\Gamma(\frac14+\frac{i\tau}2)$ \cite{BhaduriKhareLaw}, so
\begin{equation}\label{eq:freewinding}
 2\eta'(\tau)-2\theta'(\tau)=\log\pi\qquad\text{exactly, for every }\tau,
\end{equation}
verified to $10^{-14}$. That system is $\frac12(xp+px)$ --- the dilatation
generator --- on a ray, with continuous spectrum on all of $\R$, which is what
Proposition~\ref{prop:nowinding} requires. Two consequences, in opposite
directions. The delay clause is satisfied by a \emph{free} one-dimensional
Hamiltonian, so as a test of a theory it has almost no power. And what the free
generator misses is exactly the conductor: the constant $\log\pi$, which is the
$\pi^{-s/2}$ of the completion.

That the target is an inner function is the same statement:
$e^{2i\theta(\tau)}=\pi^{-i\tau}\Gamma(\frac14+\frac{i\tau}2)/
\Gamma(\frac14-\frac{i\tau}2)$, unimodular, with poles at $\tau=i(\frac12+2n)$ and
zeros at $\tau=-i(\frac12+2n)$. Proposition~\ref{prop:flux} already says
$K_\omega$ is inner; Proposition~\ref{prop:nowinding} says no correlator is.

\subsection{The expansion, and the graded invariants}\label{ssec:invariants}

The discrimination is in the expansion. Let a candidate's reflection amplitude be
unimodular of $\Gamma$-ratio type,
\begin{equation}\label{eq:gammaratio}
 S(\tau)=e^{2ic_1\tau}\prod_j
 \left(\frac{\Gamma(a_j+i\beta_j\tau)}{\Gamma(a_j-i\beta_j\tau)}\right)^{n_j},
 \qquad a_j>0,\ \beta_j>0,\ n_j\in\R\setminus\{0\},
\end{equation}
finitely many factors. The multiplicities are real \emph{of either sign}: every
amplitude in the survey of \S\ref{ssec:survey} except the $L$-function targets
themselves has at least one $\Gamma$ in a denominator, so restricting to $n_j>0$
would remove the physics rather than normalize it. From
$\re\psi(a+iy)=\log y+\sum_{m\geq1}(-1)^{m+1}\frac{B_{2m}(a)}{2my^{2m}}$:

\begin{proposition}[The delay expansion]\label{prop:expansion}
For every $M\geq1$,
\[
 \mathcal T(\tau)=\Lambda\log\tau+c_0
 +\sum_{m=1}^{M}\frac{(-1)^{m+1}}{m}\,\frac{\Sigma_{2m}}{\tau^{2m}}
 +O\big(\tau^{-2M-2}\big),
\]
\[
 \Lambda=2\sum_jn_j\beta_j,\qquad
 \Sigma_{2m}=\sum_j\frac{n_jB_{2m}(a_j)}{\beta_j^{2m-1}},
\]
the series being asymptotic and divergent, as Stirling's is,
with $c_0=2c_1+2\sum_jn_j\beta_j\log\beta_j$. The target is
$\mathcal T_{\rm Weil}=\log\tau-\log2\pi-\frac1{24\tau^2}+O(\tau^{-4})$.
\end{proposition}

The prefactor $(-1)^{m+1}/m$ is universal --- it depends on neither the amplitude
nor $m$'s companions --- so it is the same for candidate and target and cancels in
every comparison below. It is kept out of $\Sigma_{2m}$ so that the target reads
$\Sigma_{2m}^{\rm Weil}=2^{2m-1}B_{2m}(\frac14)$; it is not optional in the
equation, and at $m=1$ it is $+1$, which is why earlier versions, which wrote only
the $\tau^{-2}$ term, did not need it.

A candidate's natural momentum is not $\tau$, and the relation is linear,
$\tau=\kappa P$; matching the leading coefficient forces $\kappa=\Lambda$. Rather
than substitute by hand for each candidate --- which is what \S\ref{ssec:liouville}
does for Liouville, and where a reader's suspicion goes first --- one may pass to
quantities in which the dictionary has already cancelled.

\begin{definition}\label{def:invariants}
$I_{2m}(S):=\Lambda^{2m-1}\Sigma_{2m}$.
\end{definition}

\begin{proposition}[The invariants are dictionary-free]\label{prop:dictfree}
$I_{2m}$ is unchanged by any linear change of momentum and by any common
rescaling $\beta_j\mapsto\lambda\beta_j$. The target is
\[
 I_{2m}^{\rm Weil}=2^{2m-1}B_{2m}(\tfrac14)
 =-\tfrac1{24},\quad\tfrac7{480},\quad-\tfrac{31}{2688},\quad\tfrac{127}{7680},\ \dots
\]
\end{proposition}

\begin{proof}
Under $\tau=\kappa P$ a delay measured in $P$ becomes
$\mathcal T(\tau)=\kappa^{-1}\mathcal T_{\rm phys}(\tau/\kappa)$, so the leading
coefficient is $\Lambda/\kappa$ and $\Sigma_{2m}\mapsto\kappa^{2m-1}\Sigma_{2m}$,
the universal prefactor going along unchanged; imposing $\kappa=\Lambda$ gives the
claim. For the
second, $\Lambda\mapsto\lambda\Lambda$ and
$\Sigma_{2m}\mapsto\lambda^{-(2m-1)}\Sigma_{2m}$. The target is the single factor
$n=1$, $a=\frac14$, $\beta=\frac12$.
\end{proof}

\begin{corollary}[$\Lambda$ is the degree, once the momentum is normalized]\label{cor:degree}
For the archimedean factor $\prod_{j=1}^d\Gamma_\R(s+\mu_j)$ of a degree-$d$
$L$-function at $s=\frac12+i\tau$ one has $a_j=\frac{1/2+\mu_j}2$,
$\beta_j=\frac12$, $n_j=1$, so $\Lambda=d$ exactly.
\end{corollary}

$\Lambda$ is emphatically \emph{not} one of the dictionary-free quantities: it
scales, $\Lambda\mapsto\lambda\Lambda$, which is the whole reason for
Definition~\ref{def:invariants}. Bulk Liouville at $b=1$ has $\Lambda=8$ in its own
variable and $\Lambda=1$ after $\beta_j\mapsto\beta_j/8$, with $I_2$ and $I_4$
unmoved. "Degree one" is therefore a statement about the normalization
$\tau=\im s$, $\beta_j=\frac12$, in which an $L$-function's archimedean factor is
canonically written; it is not a test a candidate with an unknown momentum
dictionary can be put to, and no verdict in \S\ref{ssec:survey} rests on it.

\begin{theorem}[One factor: the classification is complete]\label{thm:onefactor}
Let $S$ be \eqref{eq:gammaratio} with a single factor. Then
\[
 I_{2m}(S)=2^{2m-1}n^{2m}B_{2m}(a)\qquad\text{for every }m,
\]
\emph{independent of $\beta$}. Consequently $S$ matches the target in $I_2$ and
$I_4$ if and only if
\[
 (a,n)\in\Big\{\big(\tfrac14,\pm1\big),\ \big(\tfrac34,\pm1\big),\
 \big(\tfrac12,\pm\tfrac12\big)\Big\},
\]
and each of these then matches $I_{2m}$ for every $m$.
\end{theorem}

\begin{proof}
$\Lambda=2n\beta$ and $\Sigma_{2m}=nB_{2m}(a)\beta^{-(2m-1)}$, so the powers of
$\beta$ cancel and $I_{2m}=2^{2m-1}n^{2m}B_{2m}(a)$. Put $t=n^2>0$ and
$u=B_2(a)=a^2-a+\frac16$. Matching $I_2$ and $I_4$ reads $tu=-\frac1{48}$ and
$t^2B_4(a)=\frac7{3840}$. Since $B_4(a)=(a^2-a)^2-\frac1{30}=u^2-\frac u3-\frac1{180}$,
eliminating $t$ gives
\[
 \frac{u^2-\frac u3-\frac1{180}}{2304\,u^2}=\frac7{3840}
 \iff 576u^2+60u+1=0,
\]
whose discriminant is $3600-2304=36^2$ and whose roots are $u=-\frac1{48}$ and
$u=-\frac1{12}$. The first gives $t=1$ and $a^2-a+\frac3{16}=0$, so $n=\pm1$ and
$a\in\{\frac14,\frac34\}$; the second gives $t=\frac14$ and $(a-\frac12)^2=0$, so
$n=\pm\frac12$ and $a=\frac12$. That all six then match every $I_{2m}$ is
Remark~\ref{rem:degeneracy}.
\end{proof}

\begin{corollary}[The shift must sit near the critical line]\label{cor:signrule}
For a single factor $I_2=2n^2B_2(a)$ carries the sign of $B_2(a)$ \emph{whatever
the sign of $n$}, so matching forces
\[
 \Big|a-\tfrac12\Big|<\tfrac1{2\sqrt3}=0.288675\ldots,
\]
and no single-factor amplitude with $a\geq1$ can match at any multiplicity. For
several factors with all $n_j>0$ one has $\Lambda>0$, so the same window follows
from $\Sigma_2<0$ and at least one $a_j$ lies in it.
\end{corollary}

\begin{center}
\begin{tabular}{@{}lccc@{}}
\toprule
factor, at $s=\frac12+i\tau$ & $a$ & $B_2(a)$ & in the window\\
\midrule
$\Gamma_\R(s)$ --- $\zeta$, even $\chi$ & $\frac14$ & $-\frac1{48}$ & yes\\
$\Gamma_\R(s+1)$ --- odd $\chi$ & $\frac34$ & $-\frac1{48}$ & yes\\
$\Gamma_\C(s)$ --- a complex place & $\frac12$ & $-\frac1{12}$ & yes, deepest\\
$\Gamma(1+i\,\cdot)$ --- generic reflection amplitude & $1$ & $+\frac16$ & \textbf{no}\\
$\Gamma(\frac32+i\,\cdot)$ & $\frac32$ & $+\frac{11}{12}$ & \textbf{no}\\
\bottomrule
\end{tabular}
\end{center}

\begin{remark}[The degeneracy, and what a match can certify]\label{rem:degeneracy}
Three moves leave every $I_{2m}$ fixed. \emph{Inversion} $n\mapsto-n$, since
$\Lambda$ and $\Sigma_{2m}$ both change sign and $2m-1$ is odd --- so the test
cannot distinguish $S$ from $1/S$, and does not fix the sign of the delay.
\emph{Parity} $a\mapsto1-a$, since $B_{2m}(1-a)=B_{2m}(a)$ --- so it cannot
distinguish $\Gamma_\R(s)$ from $\Gamma_\R(s+1)$, that is $\zeta$ from a real odd
character. And \emph{Gauss multiplication}: from $B_n(2x)=2^{n-1}[B_n(x)+B_n(x+\frac12)]$
at $x=\frac14$,
\[
 B_{2m}(\tfrac12)=2^{2m}B_{2m}(\tfrac14)\qquad\text{for every }m\geq1,
\]
exactly, so $(a{=}\frac14,\beta,n)$ and $(a{=}\frac12,2\beta,\frac n2)$ are
indistinguishable in $\Lambda$ and in every $\Sigma_{2m}$. Hence the three
solutions of Theorem~\ref{thm:onefactor} are a single orbit, and the invariants
identify the archimedean factor only \emph{modulo} inversion, parity and the
Legendre--Gauss multiplication relations. They are an instrument for exclusion; a
match certifies membership of an orbit, not of a point.
\end{remark}

\begin{proposition}[Rigidity at a common shift]\label{prop:rigidity}
Suppose all $a_j=a$ and all $n_j>0$. Then $I_2=-\frac1{24}$ forces
\[
 \sum_jn_j\ \leq\ \big(48|B_2(a)|\big)^{-1/2},
\]
with equality if and only if all $\beta_j$ agree. The bound equals one exactly
when $a\in\{\frac14,\frac34\}$ and is below one for $a$ strictly between them, so
\emph{if $a\in[\frac14,\frac34]$} then integer multiplicities force a single
factor with $n=1$ and $a\in\{\frac14,\frac34\}$: the archimedean factor of a
degree-one $L$-function, of either parity.
\end{proposition}

\begin{proof}
$I_2=2B_2(a)\big(\sum_jn_j\beta_j\big)\big(\sum_jn_j/\beta_j\big)$, and with
$n_j>0$ Cauchy--Schwarz gives
\[
 \Big(\sum_jn_j\Big)^2\leq\Big(\sum_jn_j\beta_j\Big)\Big(\sum_jn_j/\beta_j\Big)
 =\frac{|I_2|}{2|B_2(a)|},
\]
with equality if and only if the $\beta_j$ agree; and $|B_2(a)|=\frac1{48}$ has
roots $a=\frac14$ and $a=\frac34$.
\end{proof}

Three hypotheses here are not decoration, and earlier versions of this section
carried the statement without the third.

\emph{The restriction $a\in[\frac14,\frac34]$ is needed.} Outside it and inside
the sign window the bound exceeds one and $I_2$ alone forces nothing. At
$a=\frac{4-\sqrt5}8=0.2204915\ldots$ one has $B_2(a)=-\frac1{192}$ exactly and the
bound is $2$, attained: two factors with $n_1=n_2=1$ and equal widths give
$I_2=-\frac1{24}$ exactly, for every common $\beta$. Equivalently this is the
single factor $n=2$, $a=\frac{4-\sqrt5}8$ --- an integer multiplicity, and not
$n=1$ with $a\in\{\frac14,\frac34\}$. A scan of common-shift configurations with
$n_1,n_2\leq6$ across the window finds none that also matches $I_4$ (the example
above has $I_4=-0.678$ against $\frac7{480}$), so imposing $I_4$ appears to restore
the conclusion --- but that is a different statement and it is not proved here.
What \emph{is} proved, and needs no common shift, no positivity of $n_j$ and no
integrality, is Theorem~\ref{thm:onefactor}.

\emph{The common shift is not removable}: with $a_1=\frac12$, $a_2=1$,
$n_1=n_2=1$, $\beta_1=\frac{13-\sqrt{153}}4$, $\beta_2=\frac12-\beta_1$, both
$\Lambda=1$ and $\Sigma_2=-\frac1{24}$ hold exactly, a positive $B_2$ paying for a
negative one. \emph{Nor is integrality}: $n=\pm\frac12$ at $a=\frac12$ is the
third solution of Theorem~\ref{thm:onefactor}, and it is the one the $c=1$ matrix
model realizes (\S\ref{ssec:survey}).

\subsection{Liouville, excluded}\label{ssec:liouville}

The bulk Liouville reflection amplitude is unimodular and does wind:
$\frac{d}{dP}\arg S=4Q\log P+C_L+\frac{Q}{12P^2}+O(P^{-4})$, $Q=b+b^{-1}$.
The Liouville coordinate is the radial one, so the momenta are linearly related,
$\tau=2\kappa P$, and matching the leading coefficient \emph{forces} $\kappa=2Q$.
Nothing is then left but $b$; the cosmological constant $\mu$ enters the
prefactor linearly in $P$ and so moves only $c_0$. Substituting $P=\tau/4Q$,
\[
 \mathcal T_L(\tau)=\log\tau-\log4Q+\frac{C_L}{4Q}+\frac{Q^2}{3\tau^2}+O(\tau^{-4}) .
\]
The constant can be matched, by one $\mu$ for each $b$. The $\tau^{-2}$ term
cannot: $Q\geq2$ for every $b>0$, so $Q^2/3\geq\frac43$, against the target's
$-\frac1{24}$ --- wrong sign, and never within a factor of $32$ in magnitude.
This is Corollary~\ref{cor:signrule} with $a=1$. \textbf{Liouville is excluded at
every coupling}, and so is every \emph{single-factor} amplitude built from
$\Gamma(m+i\,\cdot)$ with $m\geq1$, by Theorem~\ref{thm:onefactor}.

\begin{remark}[The single-factor hypothesis is needed]\label{rem:retraction}
Earlier versions of this section dropped it. That was wrong. With multiplicities
of both signs, take two factors with \emph{both} shifts at $a=1$,
\[
 S_r(\tau)=\frac{\Gamma(1+ir\tau)}{\Gamma(1-ir\tau)}\cdot
 \frac{\Gamma(1-i\tau)}{\Gamma(1+i\tau)},\qquad
 I_2=-\frac{(r-1)^2}{3r},
\]
so $I_2=-\frac1{24}$ at the two reciprocal roots of $8r^2-17r+8=0$, namely
$r=\frac{17\pm\sqrt{33}}{16}$, $r=1.42153516\ldots$ --- the target's $I_2$
exactly, from two shifts in the excluded region. Using $r^2=(17r-8)/8$,
\[
 I_4=\frac{r^2+r+1}{240\,r}=\frac5{384}
 \qquad\text{against the target }\frac7{480},
\]
short by $\frac3{28}$. More is true: with the multiplicities free, the pair
$I_2=-\frac1{24}$, $I_4=\frac7{480}$ has a one-parameter family of solutions with
both shifts at $a=1$, parametrized by $P=\Lambda/2\in(0,\frac{\sqrt7}4)$ through
$\beta^{-2}=(7-16P^2)/\big(16P^2(1+8P^2)\big)$. Imposing $I_6$ as well leaves
exactly two points of that parametrization, $P=\frac14$ and $P=\frac12$ --- and
they are one amplitude, since $P$ double-covers the family modulo the width
rescaling of Proposition~\ref{prop:dictfree}, folding at
$P^\ast=0.3743862\ldots$, the root of $128P^4+32P^2-7=0$. That amplitude is the
target after Legendre duplication: at $P=\frac12$ it is
$(\frac12,1,2)$ with $(-\frac12,1,1)$, and
$\Gamma(1+2i\tau)/\Gamma(1+i\tau)=2^{2i\tau}\pi^{-1/2}\Gamma(\frac12+i\tau)$
identifies it with $(\frac12,\frac12,1)$ up to a linear phase, which is
Remark~\ref{rem:degeneracy}'s third move. So a truncation of the invariant
sequence is defeatable and the sequence itself, on this family, is not; whether
$\{I_{2m}\}_{m\geq1}$ determines an amplitude up to the degeneracy group in
general is open. Theorem~\ref{thm:onefactor} settles it for one factor.
\end{remark}

\subsection{The survey, and what the test cannot see}\label{ssec:survey}

Theorem~\ref{thm:onefactor} makes the test a classification rather than a window,
and it is two lines of arithmetic per candidate. Each row below lists the triples
$(n_j,a_j,\beta_j)$ read off the amplitude \emph{as quoted in the literature} ---
the $\Gamma$-structure of every row is a reading, not a re-derivation here --- and
then evaluates $\Lambda$ and the first two invariants.

\begin{center}\small
\setlength{\tabcolsep}{5pt}
\begin{tabular}{@{}llrrl@{}}
\toprule
amplitude & $(n_j,a_j,\beta_j)$ & $I_2$ & $I_4$ & verdict\\
\midrule
target $\Gamma_\R(s)$, even $\chi$ & $(1,\frac14,\frac12)$ & $-0.0416667$ & $0.0145833$ & ---\\
$\Gamma_\R(s+1)$, odd $\chi$ & $(1,\frac34,\frac12)$ & $-0.0416667$ & $0.0145833$ & match (parity)\\
inverted oscillator, half line, even & $(1,\frac14,\frac12)$ & $-0.0416667$ & $0.0145833$ & \textbf{match}\\
inverted oscillator, half line, odd & $(1,\frac34,\frac12)$ & $-0.0416667$ & $0.0145833$ & \textbf{match}\\
$c=1$ matrix model / $2$d string & $(\frac12,\frac12,1)$ & $-0.0416667$ & $0.0145833$ & \textbf{match}\\
duplication of the target & $(1,\frac18,\frac14),(1,\frac58,\frac14)$ & $-0.0416667$ & $0.0145833$ & match (Gauss)\\
inverted oscillator, \emph{full} line & $(1,\frac12,1)$ & $-0.166667$ & $0.233333$ & fails $I_2$, $\times4$\\
modular surface, arch.\ factor & $(-1,\frac12,\frac12)$ & $-0.166667$ & $0.233333$ & fails $I_2$, $\times4$\\
JT / Liouville-QM wall & $(1,1,2)$ & $+0.333333$ & $-0.266667$ & fails $I_2$, sign\\
$H_3^+$ ($m$-basis), level $k$ & $(-1,1,\frac2{k-2})$ & $+0.333333$ & $-0.266667$ & fails $I_2$, sign\\
bulk Liouville, $b=1$ & $(1,1,2b),(1,1,\frac2b)$ & $+1.33333$ & $-4.26667$ & fails $I_2$, $\times32$\\
FZZT boundary, $b=1$ & $(\frac12,1,2b),(\frac12,1,\frac2b)$ & $+0.333333$ & $-0.266667$ & fails $I_2$, sign\\
cigar, $k=\frac{14}3$, $n{=}1$ & four factors$^\dagger$ & $-0.0416667$ & $-0.0557292$ & $I_2$ exact, fails $I_4$\\
cigar, $k=\frac{782}3$, $n{=}7$ & four factors$^\dagger$ & $-0.0416667$ & $-0.267731$ & $I_2$ exact, fails $I_4$\\
\bottomrule
\end{tabular}
\end{center}

\noindent $^\dagger$ The cigar rows, both at $w=0$, are
$(-1,1,\frac2{k-2})$, $(-1,1,2)$ and $(1,\frac12+\frac{|n|}2,1)$ twice, so the
shifts are $1,1,1,1$ at $n=1$ and $1,1,4,4$ at $n=7$. $\Lambda=2\sum_jn_j\beta_j$ is not tabulated:
it is recoverable from the triples, and by Corollary~\ref{cor:degree} it is not
invariant under a common rescaling of the widths, so its magnitude is a property
of the quoted normalization rather than of the amplitude. Its \emph{sign} likewise
records only whether one writes $S$ or $S^{-1}$, and by
Remark~\ref{rem:degeneracy} no invariant and no verdict depends on that; the
$n_j$ printed above are pinned for the half-line rows by \eqref{eq:freewinding}
and for the modular surface by Remark~\ref{rem:modular}, and taken as quoted
otherwise. $I_2$ and $I_4$ are the columns that carry meaning.

Liouville, FZZT, ZZ, JT and $H_3^+$ all reduce to $\Gamma(1+i\,\cdot)$ and fail on
sign before any magnitude is compared; the coupling, the level and $\mu$ move
nothing that matters. The cigar is the interesting failure, and it is interesting
for the reason Remark~\ref{rem:retraction} isolates rather than the one it might
seem: at $w=0$ its shifts are \emph{fixed}, at $a=1$ for $n=1$ and at $a=4$ for
$n=7$, all of them outside the sign window, and what is tuned continuously to hit
$I_2=-\frac1{24}$ is the level $k$, entering as the \emph{width} $\frac2{k-2}$ of
a factor with $n<0$. The $k=\frac{14}3$ row therefore has all four shifts at
$a=1$: it is a member of the same phenomenon as the counterexample, and like it it
fails at $I_4$.

The single survivor is the inverted harmonic oscillator on a \emph{half} line,
whose parabolic-cylinder index is $\frac14$ in the even sector and $\frac34$ in
the odd, against $\frac12$ on the full line; the even/odd split maps onto the
parity split $\Gamma_\R(s)$ / $\Gamma_\R(s+1)$ of even and odd characters exactly
as Remark~\ref{rem:degeneracy} requires. That match is \cite{BhaduriKhareLaw},
already the motivating example of \S\ref{ssec:winds}. What
Theorem~\ref{thm:onefactor} adds is that there is nothing else: the $c=1$ matrix
model row is its Gauss-duplication image, not an independent candidate, as it must
be, the free-fermion and half-line problems being the same problem.

\begin{remark}[The test cannot see the zeros]\label{rem:modular}
The modular surface row deserves its own statement, because it is the one entry
where the Riemann zeros are literally present. For $\Gamma=PSL(2,\Z)$ the
Eisenstein scattering matrix \cite{Iwaniec} is
\[
 \varphi(s)=\frac{\Lambda_\zeta(2s-1)}{\Lambda_\zeta(2s)}
 =\sqrt\pi\,\frac{\Gamma(s-\frac12)\,\zeta(2s-1)}{\Gamma(s)\,\zeta(2s)},
 \qquad \Lambda_\zeta(u)=\pi^{-u/2}\Gamma(\tfrac u2)\zeta(u),
\]
which is \emph{not} $\xi(2s-1)/\xi(2s)$ for the $\xi$ of \eqref{eq:xi}: those
differ by the unimodular factor $\frac{2s-2}{2s}$, whose own delay is $-2\tau^{-2}$
and would move $I_2$ from $-\frac16$ to $+\frac{11}6$. The poles of $\varphi$ are
the zeros of $\Lambda_\zeta(2s)$, at $s=\rho/2$: the resonances \emph{are} the
Riemann zeros, on $\re s=\frac14$ under RH and in $0<\re s<\frac12$
unconditionally. If any system in this survey ought to pass an archimedean test
calibrated on $\zeta$, it is this one. It fails by a factor of four.

$\varphi$ is not of $\Gamma$-ratio type, so \eqref{eq:gammaratio} does not apply to
it whole; it factors as archimedean times arithmetic, and the invariants are
defined on the first. On $\re s=\frac12$ the functional equation and reality give
$\Lambda_\zeta(2s-1)=\overline{\Lambda_\zeta(2s)}$, so $\varphi$ is unimodular with
$\arg\varphi(\frac12+ir)=-2\arg\Lambda_\zeta(1+2ir)$, whose $\Gamma$ content is
$\Gamma(\frac12+ir)^{-1}$ against its conjugate: shift $\frac12$, not $\frac14$,
and $n=-1$. Hence in $\tau=2r$ the archimedean part of the delay is
\[
 \mathcal T_{\rm arch}(\tau)=-\log\frac\tau{2\pi}+\frac1{6\tau^2}+O(\tau^{-4}),
 \qquad \Lambda=-1,\qquad I_2=-\tfrac16 .
\]
The arithmetic part adds $2\re\frac{\zeta'}{\zeta}(1+i\tau)$, which is not small
--- it is $-2\sum_n\Lambda(n)n^{-1}\cos(\tau\log n)$, of mean zero and oscillatory,
still of size $0.90$ at $\tau=10^4$ where $\tau^{-4}=10^{-16}$ --- and which
therefore contributes to no asymptotic coefficient at all. It is the exact
counterpart of the prime term of Remark~\ref{rem:notprimes}, on the other side.

There is no paradox in a surface whose resonances are the zeros failing a test
calibrated on them. The resonances lie \emph{off} the line, while the delay is
read \emph{on} it, where $\varphi$ is governed by $\Lambda_\zeta$ at the edge
$u=1+2ir$ of the strip --- a region in which $\zeta$ has no zeros, so the delay has
no resonance jumps and the zeros enter it not at all. What is left on the line is
$\Gamma_\R(2s)$ rather than $\Gamma_\R(s)$, and that doubling is exactly the
$\frac14\to\frac12$ that costs the factor of four.

Two consequences, cutting in opposite directions. A pass is weak evidence: the
half-line inverted oscillator matches every invariant exactly and its spectrum has
nothing to do with the zeros. A failure is not evidence of absence: the modular
surface has the zeros as resonances and fails. The delay test is a measurement of
the archimedean factor on the critical line, and of nothing else.
\end{remark}

\begin{remark}[Where the primes are not]\label{rem:notprimes}
None of this expansion sees the primes. They are in the fluctuation
$S(\tau)=\pi^{-1}\arg\zeta(\frac12+i\tau)$, whose derivative does not decay ---
it is not asymptotically small --- but which oscillates with mean zero and so
contributes to no asymptotic coefficient. So the expansion tests whether a candidate has the right
\emph{archimedean} data, to any depth wanted, and says nothing about whether it
has any arithmetic --- two questions that were previously entangled in one test.
This is the same separation Section~\ref{sec:sampling} found on the other side,
where the margin turned out to be set by the density and not by the primes.
Remark~\ref{rem:modular} says the stronger thing: the expansion cannot see the
\emph{zeros} either, and exhibits the same mean-zero arithmetic term on the other
side. Neither reaches any asymptotic coefficient.
\end{remark}

\begin{remark}[Status of this section]\label{rem:delaystatus}
Propositions~\ref{prop:defectdelay}, \ref{prop:nowinding}, \ref{prop:expansion},
\ref{prop:dictfree} and \ref{prop:rigidity}, Theorem~\ref{thm:onefactor} and
their corollaries are written proofs, unconditional;
Proposition~\ref{prop:nowinding} quotes Bernstein's theorem. The Liouville
exclusion assumes the momentum dictionary is linear, which is stated where it is
used and which Section~\ref{sec:delay}'s other results do not need ---
Proposition~\ref{prop:dictfree} removes that assumption from everything after it.
Proposition~\ref{prop:rigidity}'s last clause needs $a\in[\frac14,\frac34]$; the
counterexample at $a=\frac{4-\sqrt5}8$ shows earlier versions of it, which omitted
that hypothesis, were wrong. The $(n_j,a_j,\beta_j)$ triples of
\S\ref{ssec:survey} are readings of the quoted amplitudes; the arithmetic on them
is a computation, and the cigar rows are the ones most in need of independent
re-derivation. In Remark~\ref{rem:modular} the archimedean expansion and the
unimodularity are proved and checked; the size of the arithmetic term is a
labelled computation, and $\re s=\frac14$ for the resonances is conditional on
RH. The one-parameter family and its
two $I_6$ roots in Remark~\ref{rem:retraction} are a bisection sweep, with the two
roots then identified in closed form. Every number
is verified in Appendix~\ref{app:checks}. Section 11.1's description of
\cite{OWL} is a reading with its scope in the footnote; the defect spectrum and
normalizations are quoted, not re-derived.
\end{remark}
